\SourceLeaf{268}
\begin{list}{}{\setlength{\leftmargin}{1.5em}\setlength{\itemindent}{-1.5em}\setlength{\itemsep}{0pt}\setlength{\parsep}{0pt}}
\item \(e\) ist die Exzentrizität der Ellipse, durch deren Drehung um ihre kleine Achse die Ellipsoidfläche entsteht.
\item \(t\) und \(90^0-w\) sind die Länge und Breite eines beliebigen Punkts auf dieser Fläche. Somit ist \(w\) der Winkel zwischen der Flächennormalen in diesem Punkt und der kleinen Achse.
\item \(T\) und \(90^0-U\) sind die Länge und Breite des entsprechenden Punkts auf der Kugelfläche.
\item \(i\) ist die imaginäre Einheit \(\BareRoot{-1}\).
\item \(f\) bezeichnet eine beliebig zu wählende Funktion.
\end{list}
\noindent Es werden stets natürliche Logarithmen verwendet.\par
\noindent Mit \(m\) bezeichnen wir das Vergrößerungsverhältnis. Dabei verhält sich jedes Linienelement auf der Ellipsoidfläche zum entsprechenden Linienelement auf der Kugelfläche wie 1 zu \(m\). Dieses Verhältnis hat an jeder Stelle der einen wie der anderen Fläche einen bestimmten Wert und kann sich von Stelle zu Stelle ändern.

Die einfachste Lösung erhält man, indem man die beliebige Funktion gleich ihrem Argument setzt, also
\[
f\upsilon=\upsilon
\]
Diese Art der Abbildung ist tatsächlich auch die geeignetste, wenn die \textit{ganze} Oberfläche des Ellipsoids auf die Kugelfläche übertragen werden soll. Bei geodätischen Rechnungen, für die immer nur ein verhältnismäßig kleiner Teil der Erdoberfläche in Betracht kommt, ist es jedoch, wie dort bereits bemerkt, viel vorteilhafter, der Funktion einen konstanten imaginären Anteil hinzuzufügen, also
\[
f\upsilon=\upsilon-i\log k
\]
Der Radius der Kugel und die Konstante \(k\) lassen sich dann so bestimmen, dass am mittleren Breitenkreis für das Vergrößerungsverhältnis \(m\) \(=1\) gilt und es sich bis zu einigen Grad nach Norden und Süden kaum merklich von 1 unterscheidet. Die praktische Verwendbarkeit für geodätische Rechnungen hängt vor allem davon ab, wie wenig sich dieses Verhältnis innerhalb des abgebildeten Gebiets ändert. Die Abweichung vom Wert 1 ist nämlich von zweiter Ordnung in Bezug auf
